\section{Limitations and Future Directions}
\label{sec:limitations}

Despite the security principles and lifecycle-aligned defense practices summarized in Section~\ref{sec:discussion}, current CUA security remains constrained by several fundamental limitations. We highlight key open problems and outline promising directions for future work.

\subsection{Attribution and Observability of Subjective Factors}
The 2$\times$2 framework distinguishes curiosity from malice, but in practice both are hard to observe. Curious overreach and malicious injection can produce similar surface behaviors (e.g., PII leakage); distinguishing ``the model explored too much'' from ``the model was induced to exfiltrate'' requires intent tracing and provenance---which remain limited~\cite{Jones2025CUA}. Behavioral boundaries (what counts as on-task) are underspecified. Future work should develop richer telemetry, causal tracing of intent, and standardized notions of task boundaries to better separate benign curiosity from induced malice.

\subsection{Lack of Centralized Security Oversight}
CUA frameworks and skill registries are largely decentralized. There is no single authority to audit baselines, enforce compliance, or coordinate vulnerability disclosure. This fragmentation enables supply-chain attacks, typosquatting, and inconsistent patching~\cite{Jones2025CUA,AgentSecurityBench2025}. Long-term progress requires stronger ecosystem governance, such as shared registries with security ratings, coordinated disclosure processes, and policy frameworks that balance openness with safety.

\subsection{Capability--Safety Inverse Scaling}
CUAHarm~\cite{CUAHarm2025} and PropensityBench~\cite{PropensityBench2025} report that more capable models often exhibit higher misuse risk. Under task pressure or with high-risk tools, models that appear safe in neutral settings may choose unsafe shortcuts. Safety is not static; operational pressure and tool access must be incorporated into evaluation frameworks. Future research should explore adaptive evaluation that conditions on task difficulty, tool configuration, and deployment context, as well as training schemes that explicitly penalize unsafe shortcuts under pressure.

\subsection{Cross-modal Injection Defenses}
Existing detectors focus on textual prompt injection; WAInjectBench~\cite{WAInjectBench2025} shows they fail against imperceptible pixel perturbations. Chameleon~\cite{Chameleon2025} demonstrates that visual and environmental injection require dedicated, multi-modal defenses. Integrating vision and text in a unified detection pipeline remains an open problem. Promising directions include multi-modal anomaly detection, robust perceptual anchoring, and joint training of vision--language defenses that are evaluated on realistic, lifecycle-aware benchmarks.

\subsection{Plan-time Safety Awareness and Evaluation Gaps}
LPS-BENCH~\cite{LPSBench2026} evaluates agents on long-horizon planning and finds they lack foresight for cascading risks. Safety interventions at execution time are reactive; addressing risks at plan-time could prevent harmful trajectories. Scaling plan-time reasoning to complex, multi-step tasks is technically challenging. At the same time, benchmarks such as OS-Harm~\cite{OSHarm2025}, CUAHarm~\cite{CUAHarm2025}, and LPS-BENCH only partially cover real deployment settings; they are not yet fully integrated into release gating and continuous regression. Future work should improve coverage of cross-modal and cross-environment threats, and more tightly couple benchmark results to deployment decisions and maintenance workflows.

\subsection{Multi-agent Isolation and Governance}
When CUAs interact via MCP or platforms such as Moltbook, trust boundaries blur. Prompt Infection~\cite{PromptInfection2024} shows malicious prompts can self-replicate across agents. Multi-Agent Security Tax~\cite{MultiAgentSecurityTax2025} analyzes the trade-off between security controls and collaboration. Establishing cryptographic identity~\cite{AegisProtocol2025}, verifiable policy compliance, and content governance for agent societies is essential but not yet standardized~\cite{AgenticAICybersecuritySurvey2026}. Future directions include standardized identity and attestation for agents, isolation policies for inter-agent communication, and governance models that align incentives for platforms, developers, and regulators.

